% =====================================================================
%  PROJECT REPORT - FireVLM (Dept. of CSE, Amrita School of Computing, Bengaluru)
%  Compile with XeLaTeX, twice. Keep "amrita_logo.png" in the same folder.
%  Anything in RED in square brackets is still a placeholder (\todo{...}).
%  All numbers below come from the Kaggle notebook outputs / repo digest.
% =====================================================================
\documentclass[12pt]{report}

\usepackage[letterpaper,
            top=1.38in, bottom=1.18in, left=1.58in, right=1.18in,
            headsep=0.38in, footskip=0.5in]{geometry}
\usepackage{fontspec}
\IfFontExistsTF{Times New Roman}{\setmainfont{Times New Roman}}{%
  \IfFontExistsTF{TeX Gyre Termes}{\setmainfont{TeX Gyre Termes}}{%
    \setmainfont{Liberation Serif}}}
\IfFontExistsTF{Calibri}{\newfontfamily\hffont{Calibri}}{%
  \IfFontExistsTF{Carlito}{\newfontfamily\hffont{Carlito}}{%
    \newfontfamily\hffont{Liberation Sans}}}

\usepackage{graphicx}
\usepackage{xcolor}
\usepackage{setspace}
\usepackage{fancyhdr}
\usepackage{titlesec}
\usepackage{titletoc}
\usepackage{enumitem}
\usepackage{array}
\newcolumntype{P}[1]{>{\raggedright\arraybackslash}p{#1}}
\usepackage{booktabs}
\usepackage{amsmath,amssymb}
\usepackage{tikz}
\usetikzlibrary{arrows.meta,positioning,shapes.geometric,calc}
\usepackage{pgfplots}
\pgfplotsset{compat=1.16}
\usepackage[labelsep=colon]{caption}
\usepackage[numbers,square,sort&compress]{natbib}
\usepackage[hidelinks]{hyperref}

\definecolor{wordred}{HTML}{FF0000}
\newcommand{\red}[1]{\textcolor{wordred}{#1}}
\newcommand{\todo}[1]{\textcolor{wordred}{[#1]}}

\setlength{\headheight}{14pt}
\emergencystretch=1em
\setlength{\parindent}{0pt}
\setlength{\parskip}{6pt}

% ---------------- Details --------------------------------------------
\newcommand{\projecttitle}{FireVLM: Parameter-Efficient Fine-Tuning of a Vision-Language Model for Wildfire Visual Question Answering}
\newcommand{\monthyear}{December 2026}
\newcommand{\guidename}{Dr. Rimjhim Padam Singh, Assistant Professor (Sl.Gd.), Department of Computer Science Engineering}

% ---------------- Headers and footers --------------------------------
\newif\ifinmain
\fancypagestyle{front}{%
  \fancyhf{}
  \renewcommand{\headrulewidth}{0pt}
  \fancyfoot[C]{\thepage}}
\fancypagestyle{main}{%
  \fancyhf{}
  \renewcommand{\headrulewidth}{0pt}
  \fancyhead[L]{\hffont\small\leftmark}
  \fancyhead[R]{\hffont\small\monthyear}
  \fancyfoot[L]{\hffont\small\vspace*{\baselineskip}Dept. of CSE, ASC, Bengaluru}
  \fancyfoot[R]{\hffont\small Page | \thepage\vspace*{\baselineskip}}}
% chapter-opening pages use "plain": same as "main" in the body, page number only in the front matter
\fancypagestyle{plain}{%
  \fancyhf{}
  \renewcommand{\headrulewidth}{0pt}
  \fancyhead[L]{\ifinmain\hffont\small\leftmark\fi}
  \fancyhead[R]{\ifinmain\hffont\small\monthyear\fi}
  \fancyfoot[C]{\ifinmain\else\thepage\fi}
  \fancyfoot[L]{\ifinmain\hffont\small\vspace*{\baselineskip}Dept. of CSE, ASC, Bengaluru\fi}
  \fancyfoot[R]{\ifinmain\hffont\small Page | \thepage\vspace*{\baselineskip}\fi}}
\renewcommand{\chaptermark}[1]{\markboth{#1}{}}
\renewcommand{\sectionmark}[1]{}

\makeatletter
\renewcommand{\chapter}{\global\@topnum\z@\@afterindentfalse
  \secdef\@chapter\@schapter}
\makeatother

\titleformat{\chapter}[display]
  {\normalfont\fontsize{14}{17}\bfseries\centering}
  {CHAPTER - \thechapter}{6pt}{\MakeUppercase}
\titlespacing*{\chapter}{0pt}{0pt}{18pt}
\titleformat{\section}{\normalfont\fontsize{12}{14}\bfseries}{\thesection}{0.6em}{}
\titlespacing*{\section}{0pt}{12pt}{6pt}
\titleformat{\subsection}{\normalfont\fontsize{12}{14}\bfseries}{\thesubsection}{0.6em}{}
\titlespacing*{\subsection}{0pt}{10pt}{4pt}
\setcounter{secnumdepth}{3}
\setcounter{tocdepth}{2}

\contentsmargin{0.5in}
\newcommand{\tocdots}{\titlerule*[0.6pc]{.}}
\titlecontents{chapter}[0pt]{\addvspace{3pt}}
  {CHAPTER \thecontentslabel\ -- \MakeUppercase}
  {\hspace*{0.25in}\MakeUppercase}
  {\tocdots\contentspage}
\titlecontents{section}[0.95in]{}
  {\contentslabel{0.45in}\MakeUppercase}{\MakeUppercase}
  {\tocdots\contentspage}
\titlecontents{subsection}[1.55in]{}
  {\contentslabel{0.6in}\MakeUppercase}{\MakeUppercase}
  {\tocdots\contentspage}
\titlecontents{figure}[0.55in]{\addvspace{2pt}}
  {\contentslabel{0.55in}}{}{\tocdots\contentspage}
\titlecontents{table}[0.55in]{\addvspace{2pt}}
  {\contentslabel{0.55in}}{}{\tocdots\contentspage}

\renewcommand{\contentsname}{TABLE OF CONTENTS}
\renewcommand{\listfigurename}{List Of Figures}
\renewcommand{\listtablename}{List Of Tables}
\renewcommand{\bibname}{REFERENCES}

\captionsetup[figure]{name=Fig., font=bf, labelfont=bf, skip=6pt}
\captionsetup[table]{name=Table, font=bf, labelfont=bf, position=top, skip=6pt}

\renewcommand{\bibnumfmt}[1]{[#1].}
\setlength{\bibsep}{2pt}

\newcommand{\frontchapter}[1]{%
  \clearpage\phantomsection
  \addcontentsline{toc}{chapter}{#1}%
  \chapter*{#1}}

% TikZ styles for diagrams
\tikzset{
  blk/.style={draw, rounded corners=2pt, align=center, minimum height=13mm,
              text width=25mm, font=\footnotesize, fill=gray!8},
  arr/.style={-{Stealth[length=2.2mm]}, thick}
}

\begin{document}
\addtocontents{toc}{\protect\hfill\textbf{Page No.}\protect\par}

% =====================================================================
%  TITLE PAGE
% =====================================================================
\begin{titlepage}
\centering
\vspace*{0.4in}
{\fontsize{14}{17}\selectfont\bfseries \projecttitle\par}
\vspace{10pt}
A PROJECT REPORT\par
\vspace{10pt}
{\fontsize{14}{17}\selectfont\bfseries\itshape Submitted by\par}
\vspace{14pt}
\begin{tabular}{@{}l@{\hspace{0.45in}}l@{}}
BL.EN.U4AIE23123 & Praanesh Balakrishnan Nair\\[16pt]
BL.EN.U4AIE23126 & Rishab Ramesh Nair\\[16pt]
BL.EN.U4AIE23130 & Sanjushree Rajan\\[16pt]
BL.EN.U4AIE23135 & Varun Adhitya G B\\[16pt]
\end{tabular}\par
\vspace{18pt}
{\fontsize{14}{17}\selectfont\bfseries\itshape in partial fulfillment for the award of the degree of\par}
\vspace{18pt}
{\fontsize{14}{17}\selectfont\bfseries BACHELOR OF TECHNOLOGY\par}
\vspace{16pt}
{\small IN\par}
\vspace{16pt}
COMPUTER SCIENCE AND ENGINEERING\par
(Artificial Intelligence) \\
\vspace{12pt}
\IfFileExists{amrita_logo.png}{\includegraphics[width=1.75in]{amrita_logo.png}}{\todo{amrita\_logo.png not found}}\par
\vspace{12pt}
AMRITA SCHOOL OF COMPUTING, BENGALURU\par
\vspace{12pt}
{\fontsize{14}{17}\selectfont AMRITA VISHWA VIDYAPEETHAM\par}
\vspace{12pt}
BENGALURU 560 035\par
\vspace{12pt}
\MakeUppercase{\monthyear}\par
\end{titlepage}

% =====================================================================
%  FRONT MATTER
% =====================================================================
\pagestyle{front}
\pagenumbering{roman}
\setcounter{page}{1}

\frontchapter{ACKNOWLEDGEMENTS}
\begin{spacing}{1.5}
\setlength{\parindent}{0.5in}
\indent The satisfaction that accompanies successful completion of any task would
be incomplete without mention of people who made it possible, and whose
constant encouragement and guidance have been source of inspiration
throughout the course of this project work.

We offer our sincere pranams at the lotus feet of \textbf{``AMMA'',
MATA AMRITANANDAMAYI DEVI} who showered her blessing upon us throughout the
course of this project work.

We owe our gratitude to \textbf{Prof. Manoj P.}, Director, Amrita Vishwa Vidyapeetham Bengaluru Campus. We would like to place our heartfelt gratitude to \textbf{Dr. Gopalakrishnan E.A.,} Principal, Amrita School of Computing and Amrita School of Artificial Intelligence, Bengaluru for his
valuable support and inspiration.

It is a great pleasure to express our gratitude and indebtedness to our
project guide \textbf{\guidename}, Amrita School of Computing, Bengaluru for
her valuable guidance, encouragement, moral support, and affection
throughout the project work.

We would like to express our gratitude to project panel members for
their suggestions, encouragement, and moral support during the process of
project work and all faculty members for their academic support. Finally,
we are forever grateful to our parents, who have loved, supported and
encouraged us in all our endeavors.
\end{spacing}

\frontchapter{ABSTRACT}
\begin{spacing}{1.5}
Early and accurate situational awareness is critical in wildfire
response, yet interpreting aerial imagery of an active fire still
depends on trained analysts. This project presents \textbf{FireVLM}, a
vision-language model that answers natural-language questions about
paired RGB and thermal wildfire imagery. We use the WildFireVQA
question--answer set built on the FLAME~3 Sycan Marsh imagery (738 image
pairs, 34 questions each, 25,092 question--answer pairs) and compare two
ways of adapting a language model to this task. The first is a frozen
Qwen2.5-VL-3B-Instruct model prompted with zero to ten solved examples.
The second is a compact vision-language model that combines a frozen
CLIP ViT-B/32 encoder, trainable RGB and thermal projectors, and a
4-bit quantised 1.5B-parameter Qwen2 language model fine-tuned with
QLoRA in several adapter configurations (query/value; query/key/value/output;
all linear layers; ranks 4 to 16). All runs use free Kaggle T4 GPUs. On the
held-out test split the frozen few-shot baseline reaches 15.0\% exact match
with no examples and 28.0\% with ten examples (200 test queries), whereas
the QLoRA models reach 84.1--86.6\% exact match on 2,482 test queries while
training only 5.8--14.0 million parameters. The best Qwen2 configuration
(query/key/value/output adapters, rank 4) reaches 86.62\% exact match, and
larger adapters did not help: rank 16 and all-linear adapters scored lower. A
Llama-3.2-1B backbone trained with the same pipeline reached up to 88.48\%.
The results show that a small adapter budget is sufficient for fire-scene
question answering and give a practical recipe for low-resource adaptation of
vision-language models to disaster imagery.
\end{spacing}

\clearpage
{\linespread{0.95}\selectfont\tableofcontents}

\clearpage\phantomsection
\addcontentsline{toc}{chapter}{LIST OF FIGURES}
\listoffigures

\clearpage\phantomsection
\addcontentsline{toc}{chapter}{LIST OF TABLES}
\listoftables

% =====================================================================
%  MAIN CHAPTERS
% =====================================================================
\clearpage
\pagestyle{main}
\inmaintrue
\pagenumbering{arabic}
\setcounter{page}{1}
\onehalfspacing

% ---------------- CHAPTER 1 ------------------------------------------
\chapter{Introduction}

\section{Introduction to Wildfire Visual Question Answering}
Wildfires cause large losses of life, property and ecosystems, and their
frequency and intensity are increasing. Aerial platforms such as drones
now capture large volumes of visible and thermal imagery of fires, but
turning that imagery into decisions (Is there active flame? Where is the
smoke heading? Which area is burned?) is still largely manual.
Visual Question Answering (VQA) lets an operator ask such questions in
plain language and receive an answer grounded in the image
\cite{antol2015}. Recent vision-language models (VLMs) such as
Qwen-VL \cite{bai2023} are strong general-purpose VQA systems, but they
are not trained on fire scenes and are too large to fully fine-tune on
modest hardware.

This project, \textbf{FireVLM}, studies how to adapt a language model to
wildfire VQA cheaply. We use the WildFireVQA question--answer set built
over the FLAME~3 Sycan Marsh imagery \cite{wildfirevqa,flame3} and compare a frozen
few-shot baseline (Qwen2.5-VL-3B-Instruct \cite{bai2025}) with QLoRA
\cite{dettmers2023} fine-tuned models that combine a frozen CLIP vision
encoder \cite{radford2021}, trainable projectors and a 4-bit Qwen2 language
model \cite{yang2024}. The task takes a paired RGB and thermal image and a
question as input and produces a short textual answer.

\section{Motivation}
Emergency agencies need quick, consistent reading of fire imagery, and
a model that can answer free-form questions is more flexible than a
fixed classifier. Fully fine-tuning a multi-billion-parameter VLM needs
more GPU memory than a student project or a small agency can afford.
Parameter-efficient fine-tuning with low-rank adapters and 4-bit
quantisation reduces this cost sharply, but it is unclear how much
adapter capacity a fire-scene task actually needs. This project answers
that with a controlled comparison of adapter placement and rank.

\section{Problem Statement}
Given a paired aerial RGB and thermal wildfire image and a natural-language question,
produce the correct short answer, while keeping the number of trainable
parameters and GPU memory low enough to run on free-tier Kaggle T4 GPUs.

\section{Objectives}
\begin{itemize}[leftmargin=0.5in, itemsep=2pt]
  \item Establish a frozen pretrained VLM with 0--10-shot prompting as a baseline on WildFireVQA.
  \item Fine-tune a compact CLIP + projector + 4-bit language-model VLM with QLoRA in several adapter configurations (query/value, query/key/value/output, all linear layers).
  \item Compare accuracy against trainable-parameter count and training cost.
  \item Identify the smallest adapter configuration that gives most of the gain.
\end{itemize}

\section{Organisation of the Report}
Chapter~2 surveys related work. Chapter~3 lists requirements, Chapter~4
gives the design and Chapter~5 the implementation. Chapters~6 and~7 cover
testing and results, and Chapter~8 concludes with future scope.

% ---------------- CHAPTER 2 ------------------------------------------
\clearpage
\chapter{Literature Survey}

\section{Visual Question Answering}
VQA was introduced as a task where a system answers free-form questions
about an image \cite{antol2015}. Modern systems connect a vision encoder
to a large language model so that image features are consumed as tokens
by the language model.

\section{Qwen and Llama Language Models}
The Qwen-VL family couples a vision encoder with the Qwen language model
and is trained on large image--text corpora \cite{bai2023}. Its
successors improve dynamic-resolution image handling \cite{wang2024,bai2025}.
Two Qwen checkpoints are used in this project. The frozen few-shot baseline uses
\texttt{Qwen/Qwen2.5-VL-3B-Instruct}, a native vision-language model. The fine-tuned models use
\texttt{mit-han-lab/Qwen2-VL-1.5B-Instruct}~\cite{mithan}. According to its model card and
configuration, this repository contains only the extracted Qwen2 language
component (\texttt{Qwen2ForCausalLM}) without a vision tower, so we pair
it with our own CLIP encoder and projectors instead of loading it as a native
Qwen2-VL model. For a backbone comparison we also use
\texttt{meta-llama/Llama-3.2-1B-Instruct}~\cite{llama32} with the same pipeline.

\section{Vision Encoder and Multimodal Projection}
CLIP \cite{radford2021} provides image features aligned with text. We use the
frozen ViT-B/32 variant, which turns a $224\times224$ image into
$7\times7=49$ patch tokens of dimension 768. Small multilayer-perceptron
projectors map these tokens into the embedding space of the language model, a
common design in visual instruction-tuning systems.

\section{Few-Shot In-Context Learning}
Large models can adapt to a new task from a handful of examples placed
in the prompt, without weight updates \cite{brown2020}. This is the
setting of our frozen baseline notebook, where $k\in\{0,\dots,10\}$ solved
examples precede the test question.

\section{LoRA and QLoRA Fine-Tuning}
LoRA freezes the pretrained weights and learns a low-rank update
$\Delta W = BA$ for selected weight matrices \cite{hu2022}. For a layer
$W \in \mathbb{R}^{d\times k}$ the adapter adds $r(d+k)$ trainable
parameters instead of $dk$, where $r$ is the rank. QLoRA stores the
frozen base model in 4-bit precision and back-propagates through it into
the adapters, which makes fine-tuning large models possible on a single
modest GPU \cite{dettmers2023}.

\section{Wildfire Imagery Datasets}
The FLAME family provides aerial imagery of prescribed burns captured
with visible and thermal cameras. FLAME~3 adds radiometric thermal TIFF
files with per-pixel temperatures \cite{flame3}. WildFireVQA is built on FLAME~3
and contains 6,097 RGB--thermal samples from three prescribed burns (Sycan Marsh,
Willamette and Shoetank), each paired with 34 questions, for 207,298 questions
in total \cite{wildfirevqa}. The questions (6 binary and 28 multiple-choice) fall into six
categories: presence and detection, classification, distribution and
segmentation, localization and direction, cross-modal reasoning, and flight
planning. Answers were generated with Gemini~2.5~Pro from the RGB image, the
colour-mapped thermal image and a radiometric temperature summary, and were then
corrected using deterministic sensor-based rules, manual verification and
intra-frame and inter-frame consistency checks \cite{wildfirevqa}. The WildFireVQA
paper evaluates frozen multimodal models zero-shot in a multiple-choice format;
the best result is 54.76\% accuracy (Qwen3-VL-8B with RGB input and retrieved
thermal statistics) \cite{wildfirevqa}.

\begin{table}[htbp]
  \centering
  \caption{Summary of related approaches and the gap addressed}
  \label{tab:related}
  \begin{tabular}{P{1.3in}P{1.7in}P{1.8in}}
    \toprule
    \textbf{Approach} & \textbf{Strength} & \textbf{Gap for wildfire VQA} \\
    \midrule
    Few-shot prompting \cite{brown2020} & No training needed & Limited domain adaptation \\
    LoRA \cite{hu2022} & Few trainable parameters & Adapter placement not studied on fire scenes \\
    QLoRA \cite{dettmers2023} & Fits on small GPUs & Capacity versus accuracy unclear for this domain \\
    Zero-shot MLLMs on WildFireVQA \cite{wildfirevqa} & Strong general baselines & No fine-tuned, low-cost models reported \\
    \bottomrule
  \end{tabular}
\end{table}

% ---------------- CHAPTER 3 ------------------------------------------
\clearpage
\chapter{System Requirements and Analysis}

\section{Software Requirements}
\begin{itemize}[leftmargin=0.5in, itemsep=2pt]
  \item Kaggle Notebooks (Python 3.13) as the execution environment, with PyTorch 2.11.0 (CUDA 12.8)
  \item Hugging Face \texttt{transformers} 4.57.1, \texttt{accelerate} 1.10.1, \texttt{peft} 0.17.1 and \texttt{bitsandbytes} 0.48.1 (versions pinned in the Qwen2 fine-tuning notebooks)
  \item \texttt{qwen-vl-utils} for the Qwen2.5-VL few-shot notebook (installed together with \texttt{transformers} 4.49 or newer; exact versions not pinned)
  \item \texttt{kagglehub} to download the WildFireVQA and FLAME~3 datasets, plus pandas, matplotlib, Pillow and tqdm
\end{itemize}

\section{Hardware Requirements}
Training and evaluation were run on Kaggle notebooks with two NVIDIA Tesla T4
GPUs (14.6\,GiB each) and \todo{system RAM}. The fine-tuning notebooks place the model
on a single GPU (cuda:0); the second GPU is intentionally left unused. The
few-shot notebook also runs on one GPU. No local hardware beyond a web browser is needed.

\section{Functional and Non-Functional Requirements}
The system must accept a paired RGB and thermal image and a question and return a textual
answer, and must report accuracy on a held-out split. It should fit within
the memory of a 14.6\,GiB T4 GPU and finish within the session limits of the
platform; the fine-tuning notebooks therefore cache CLIP features once, use
FP16 and 4-bit weights, and keep the effective batch size at 16 through
gradient accumulation.

% ---------------- CHAPTER 4 ------------------------------------------
\clearpage
\chapter{System Design}

\section{High Level Design}
Figure~\ref{fig:pipeline} shows the fine-tuned architecture. WildFireVQA
samples are paired with their FLAME~3 RGB and thermal images, and both images are resized to
$224\times224$. A frozen CLIP ViT-B/32 encoder converts each image into 49 patch
tokens (98 visual tokens per sample). Two separate projectors, one for RGB
and one for thermal, map the tokens to the language-model width, and learned
modality (type) embeddings are added. The visual tokens are inserted into a chat-formatted
prompt together with the question, and the 4-bit language model with LoRA adapters
generates the answer, which is compared with the reference. The frozen few-shot baseline instead
feeds the RGB and thermal images of each demonstration and of the query
directly to Qwen2.5-VL-3B-Instruct as a multi-turn chat, with no training.

\begin{figure}[htbp]
  \centering
  \resizebox{\textwidth}{!}{%
  \begin{tikzpicture}[node distance=6mm and 8mm]
    \node[blk] (rgb) {RGB image\\$224\times224$};
    \node[blk, below=of rgb] (th) {Thermal image\\$224\times224$};
    \node[blk, right=of rgb] (c1) {CLIP ViT-B/32\\(frozen)};
    \node[blk, right=of th] (c2) {CLIP ViT-B/32\\(frozen)};
    \node[blk, right=of c1] (p1) {RGB projector\\(MLP, trained)};
    \node[blk, right=of c2] (p2) {Thermal projector\\(MLP, trained)};
    \node[blk, text width=28mm, fill=orange!15] (llm) at ($(p1)!0.5!(p2)+(38mm,0)$) {4-bit language model\\+ LoRA adapters};
    \node[blk, above=of llm] (q) {Question\\(chat prompt)};
    \node[blk, right=of llm] (ans) {Answer text\\(max.\ 32 tokens)};
    \draw[arr] (rgb) -- (c1);
    \draw[arr] (th) -- (c2);
    \draw[arr] (c1) -- (p1);
    \draw[arr] (c2) -- (p2);
    \draw[arr] (p1) -- (llm);
    \draw[arr] (p2) -- (llm);
    \draw[arr] (q) -- (llm);
    \draw[arr] (llm) -- (ans);
  \end{tikzpicture}}
  \caption{Architecture of the fine-tuned FireVLM models ($d$ = hidden width of the language model: 1536 for Qwen2-1.5B, 2048 for Llama-3.2-1B).}
  \label{fig:pipeline}
\end{figure}

\section{Low Level Design}
\subsection{Dataset and Data Splits}
The Kaggle release of WildFireVQA contains 207,306 annotation rows over
three burns. We keep the Sycan Marsh rows (25,092) and match every row to the
FLAME~3 Sycan Marsh Kaggle subset, which contains 738 RGB--thermal pairs (622 fire and 116
no-fire). A strict pairing audit checks that each question is attached to the right image
pair (RGB and thermal files differ, one pair per source image, 34 questions
per image). The pairs are then split by image with seed 42 into train, validation and
test sets (Table~\ref{tab:split}), so no image appears in two splits. The prompt contains only the
question; the answer options are not given to the model, so answers are
generated as free-form short text and scored by string match.

\begin{table}[htbp]
  \centering
  \caption{Image-grouped split used by the fine-tuning notebooks}
  \label{tab:split}
  \begin{tabular}{lccc}
    \toprule
    \textbf{Split} & \textbf{Image pairs} & \textbf{QA rows} & \textbf{Share} \\
    \midrule
    Train & 592 & 20,128 & 80\% \\
    Validation & 73 & 2,482 & 10\% \\
    Test & 73 & 2,482 & 10\% \\
    \bottomrule
  \end{tabular}
\end{table}

\subsection{Experiment Configurations}
Table~\ref{tab:variants} lists the configurations. Adapters are
attached to the attention projections (query $q$, key $k$, value $v$,
output $o$) or to all linear layers of the language model (attention plus the MLP projections
gate, up and down). Additional ranks were run to study capacity
(Chapter~7). Every fine-tuned run also trains the two projectors and the modality embeddings.

\begin{table}[htbp]
  \centering
  \caption{The experimental configurations (notebook names are listed in Chapter~5)}
  \label{tab:variants}
  \begin{tabular}{P{1.7in}P{1.5in}cc}
    \toprule
    \textbf{Variant} & \textbf{Adapted modules} & \textbf{Rank $r$} & \textbf{$\alpha$} \\
    \midrule
    Few-shot, frozen Qwen2.5-VL-3B & none & --- & --- \\
    LoRA efficient & $q,v$ & 8 (also 16) & $2r$ \\
    LoRA balanced & $q,k,v,o$ & 16 (also 4, 8) & $2r$ \\
    LoRA high capacity & all linear & 16 (also 8) & $2r$ \\
    \bottomrule
  \end{tabular}
\end{table}

\subsection{Hyperparameters}
\begin{table}[htbp]
  \centering
  \caption{Training hyperparameters of the Qwen2-1.5B fine-tuning runs}
  \label{tab:hparams}
  \small
  \begin{tabular}{P{2.3in}P{2.5in}}
    \toprule
    \textbf{Parameter} & \textbf{Value} \\
    \midrule
    LoRA alpha / dropout & $2r$ / 0.05 \\
    Optimiser / learning rate & AdamW / $2\times10^{-4}$, weight decay 0.01 \\
    Scheduler & Cosine with 5\% warm-up, gradient clipping 1.0 \\
    Batch size / grad.\ accumulation & 4 / 4 (effective batch 16) \\
    Epochs & Up to 100, early stopping on validation loss (patience 3) \\
    Quantisation & 4-bit NF4, double quantisation, FP16 compute \\
    Visual input & CLIP ViT-B/32, $224\times224$, 98 visual tokens per sample \\
    Projector & $768\to1024\to d$ with GELU, FP32 \\
    Max.\ answer length & 32 tokens (training target and generation) \\
    Decoding & Greedy (no sampling, one beam) \\
    Train / validation / test split & 20,128 / 2,482 / 2,482 rows (Table~\ref{tab:split}) \\
    Random seed & 42 \\
    \bottomrule
  \end{tabular}
\end{table}

The Llama-3.2-1B runs use the same pipeline but not identical settings: for example, the
$q,k,v,o$ rank-16 run used a learning rate of $1.5\times10^{-4}$, at most 8 epochs and early-stopping
patience 2 (tolerance $10^{-3}$), whereas the $q,v$ rank-8 run used $2\times10^{-4}$ and patience 3 as in
Table~\ref{tab:hparams}. The few-shot baseline uses the same
4-bit NF4 loading (double quantisation, FP16), images of $224\times224$ pixels, greedy decoding
and at most 32 new tokens.

\subsection{Trainable Parameter Count}
For one adapted matrix of size $d\times k$ and rank $r$, the number of
trainable parameters is $r(d+k)$. For Qwen2-1.5B (28 layers, hidden width 1536, 2 key/value
heads of size 128, MLP width 8960) this gives, per layer, $r(3072)$ for $q$, $r(1792)$ for each
of $k$ and $v$, $r(3072)$ for $o$ and $r(10496)$ for each of the gate, up and down
projections. The two projectors and the two modality embeddings add 4,726,784 trainable parameters
in every fine-tuned run. Table~\ref{tab:params} lists the totals; the values for the runs
that were executed match the counts printed by the notebooks.

\begin{table}[htbp]
  \centering
  \caption{Trainable parameters of the Qwen2-1.5B variants}
  \label{tab:params}
  \small
  \begin{tabular}{lccc}
    \toprule
    \textbf{Variant} & \textbf{LoRA} & \textbf{Projectors + type emb.} & \textbf{Total trainable} \\
    \midrule
    $q,v$, $r{=}8$ & 1,089,536 & 4,726,784 & 5,816,320 \\
    $q,k,v,o$, $r{=}16$ & 4,358,144 & 4,726,784 & 9,084,928 \\
    all-linear, $r{=}16$ & 18,464,768 & 4,726,784 & 23,191,552$^{*}$ \\
    \bottomrule
  \end{tabular}\\[2pt]
  {\footnotesize $^{*}$Computed from the formula above; this notebook has no saved output, so the value was not printed by a run.}
\end{table}

% ---------------- CHAPTER 5 ------------------------------------------
\clearpage
\chapter{System Implementation}

\section{Modules Used with Description}
\begin{itemize}[leftmargin=0.5in, itemsep=2pt]
  \item \textbf{Data module}: downloads WildFireVQA and FLAME~3 Sycan Marsh with \texttt{kagglehub}, runs the pairing audit and builds the image-grouped splits.
  \item \textbf{Vision module}: encodes every unique image once with the frozen CLIP ViT-B/32 and caches the 49 patch features, so CLIP is not run during training.
  \item \textbf{Prompt module}: converts each sample into the chat template of the language model, with the RGB and thermal visual tokens placed before the question; for the baseline it builds multi-turn chats with $k$ solved demonstrations.
  \item \textbf{Model module}: loads the language model in 4-bit NF4 precision, attaches LoRA adapters with \texttt{peft}, and adds the RGB and thermal projectors and modality embeddings.
  \item \textbf{Training module}: supervised fine-tuning with a causal language-modelling loss on the answer tokens only; the best checkpoint is chosen by validation loss.
  \item \textbf{Evaluation module}: generates answers greedily and computes the metrics below.
\end{itemize}

\section{Implementation Details}
Each configuration is implemented as a separate Kaggle notebook so that
runs are independent and reproducible. The main four notebooks are:
\begin{itemize}[leftmargin=0.5in, itemsep=2pt]
  \item \texttt{firevlm-few-shot-qwen-pre-trained}
  \item \texttt{firevlm-qwen2-vl-1.5b-lora-efficient-qv-r8}
  \item \texttt{firevlm-qwen2-vl-1.5b-lora-balanced-qkvo-r16}
  \item \texttt{firevlm-qwen2-vl-1.5b-lora-high-capacity-all-linear-r16}
\end{itemize}
Further notebooks with the same names but ranks 4, 8 and 16 (and the \texttt{firevlm-lora-*}
notebooks for the Llama-3.2-1B backbone) supply the capacity ablation in Chapter~7.

\subsection{Few-Shot Baseline}
The base model, Qwen2.5-VL-3B-Instruct, is loaded in 4-bit and kept frozen
(the notebook asserts that the number of trainable parameters is zero). The
test queries are 200 rows sampled with seed 42 from the images held out as the test set
(10\% of the images); examples are drawn from the remaining images. For each query a
deterministic ordered pool of 10 solved examples (from the same question category where
possible) is built, and the $k$-shot run uses the first $k$ of them, so $k$ is the
only factor that changes: $k=0,1,\dots,10$. Each example contributes its RGB
and thermal images, the question and the reference answer. Confidence intervals
use 1,000 bootstrap resamples.

\subsection{QLoRA Variants}
The base weights are quantised to 4 bits and frozen. Only the adapter
matrices, the two projectors and the modality embeddings are trained. The variants differ in the
modules that receive adapters and in rank, as in Table~\ref{tab:variants}.
Validation loss is computed each epoch; the best projectors and adapters are saved,
and the saved best checkpoint is used for the final test generation.

\subsection{Answer Generation and Scoring}
Answers are generated greedily (one beam, no sampling, at most 32 new tokens);
only the first line of the output is kept. A prediction counts as an exact match (EM) when it
equals the reference after lower-casing and trimming white space. The notebooks also report
a smoothed sentence BLEU-4 and ROUGE-L (F-measure of the longest common subsequence).
The few-shot notebook additionally reports a normalised EM (lower-casing, removing articles
and punctuation), which was identical to EM in every run. Because most answers are one or two
words, BLEU-4 is very small and is not used for comparison; EM and ROUGE-L are the main metrics.

% ---------------- CHAPTER 6 ------------------------------------------
\clearpage
\chapter{System Testing}
Testing consisted of (i) a strict pairing audit of the data, which passed with 738 image
pairs and exactly 34 questions per image, and a check that RGB and thermal paths never
coincide; (ii) checking that only the intended weights have gradients enabled (the few-shot notebook
asserts zero trainable parameters, and the fine-tuning notebooks print and verify the adapter and projector counts);
(iii) asserting that train, validation and test image sets are disjoint (and that few-shot
demonstration images are disjoint from test images); and (iv) a memory check, in that every
configuration listed in Chapter~7 completed on a 14.6\,GiB T4 GPU. Each reported result comes from a single run per
configuration with seed 42; no repeated runs with other seeds were made.

\begin{table}[htbp]
  \centering
  \caption{Test cases}
  \label{tab:tests}
  \begin{tabular}{P{1.5in}P{2.2in}c}
    \toprule
    \textbf{Test} & \textbf{Expected behaviour} & \textbf{Status} \\
    \midrule
    Data loading and pairing & 738 pairs, 34 questions per image, RGB $\neq$ thermal & Pass \\
    Trainable parameters & Only adapters, projectors and type embeddings require gradients; zero for few-shot & Pass \\
    Split leakage & No overlap between train, validation and test images & Pass \\
    Memory & Run completes within GPU memory & Pass \\
    Repeatability & Results stable over several seeds & Not tested \\
    \bottomrule
  \end{tabular}
\end{table}

% ---------------- CHAPTER 7 ------------------------------------------
\clearpage
\chapter{Results and Analysis}

\section{Evaluation Setup}
The fine-tuned models are evaluated on the 2,482 test rows (73 held-out image pairs). The
few-shot baseline is evaluated on 200 queries sampled from its own held-out images,
because every query is run 11 times with up to 22 images in the prompt. The two test sets
are therefore not identical, and the few-shot notebook resolves image paths with a simpler
mapping that yields 480 distinct RGB files instead of the 738 pairs of the audited pipeline.
The few-shot numbers are therefore a reference point and not a strictly paired comparison.
Within the fine-tuned runs the test set, metric and decoding are identical, so differences reflect the adaptation
method (and, for the two backbones, the language model).

\section{Quantitative Results}
\begin{table}[htbp]
  \centering
  \caption{Comparison of the four configurations (exact match on the respective test sets)}
  \label{tab:results}
  \small
  \setlength{\tabcolsep}{4pt}
  \begin{tabular}{P{1.65in}ccccc}
    \toprule
    \textbf{Configuration} & \textbf{Trainable params} & \textbf{EM (\%)} & \textbf{ROUGE-L} & \textbf{Train time} \\
    \midrule
    Few-shot, frozen, $k{=}10$ (Qwen2.5-VL-3B; 200 queries) & 0 & 28.00 & 0.376 & --- \\
    LoRA $q,v$, $r{=}8$ (Qwen2-1.5B) & 5,816,320 & 86.14 & 0.8932 & 7 h 58 min \\
    LoRA $q,k,v,o$, $r{=}16$ (Qwen2-1.5B) & 9,084,928 & 84.13 & 0.8773 & 6 h 42 min \\
    LoRA all-linear, $r{=}16$ (Qwen2-1.5B) & 23,191,552$^{*}$ & \todo{not run} & \todo{not run} & \todo{not run} \\
    \bottomrule
  \end{tabular}\\[2pt]
  {\footnotesize Train time is the sum of all epochs, including the three epochs after the best checkpoint. $^{*}$Computed, see Table~\ref{tab:params}.}
\end{table}

\subsection{Few-Shot Baseline}
\begin{table}[htbp]
  \centering
  \caption{Frozen Qwen2.5-VL-3B-Instruct versus number of demonstrations $k$ (200 test queries)}
  \label{tab:shots}
  \small
  \setlength{\tabcolsep}{5pt}
  \begin{tabular}{cccccc}
    \toprule
    \textbf{$k$} & \textbf{EM (\%)} & \textbf{ROUGE-L} & \textbf{Words / answer} & \textbf{Latency (s)} & \textbf{Input tokens} \\
    \midrule
    0 & 15.00 & 0.2438 & 2.51 & 0.40 & 211 \\
    1 & 20.00 & 0.3113 & 1.73 & 0.51 & 381 \\
    2 & 23.50 & 0.3280 & 1.54 & 0.65 & 552 \\
    3 & 23.50 & 0.3290 & 1.54 & 0.80 & 723 \\
    4 & 24.50 & 0.3362 & 1.60 & 0.98 & 893 \\
    5 & 25.50 & 0.3421 & 1.66 & 1.14 & 1,064 \\
    6 & 26.00 & 0.3474 & 1.67 & 1.33 & 1,235 \\
    7 & 27.00 & 0.3585 & 1.66 & 1.54 & 1,405 \\
    8 & 26.00 & 0.3468 & 1.66 & 1.76 & 1,576 \\
    9 & 27.50 & 0.3668 & 1.64 & 1.99 & 1,747 \\
    10 & 28.00 & 0.3760 & 1.61 & 2.23 & 1,919 \\
    \bottomrule
  \end{tabular}
\end{table}

\begin{figure}[htbp]
  \centering
  \begin{tikzpicture}
    \begin{axis}[width=0.85\textwidth, height=6.2cm, xlabel={Number of demonstrations $k$},
      ylabel={Score (\%)}, xmin=0, xmax=10, ymin=10, ymax=42, xtick={0,1,...,10},
      legend pos=south east, grid=major, grid style={gray!25}, every axis plot/.append style={thick}]
      \addplot[mark=*, color=blue!70!black] coordinates {(0,15)(1,20)(2,23.5)(3,23.5)(4,24.5)(5,25.5)(6,26)(7,27)(8,26)(9,27.5)(10,28)};
      \addlegendentry{Exact match}
      \addplot[mark=square*, color=red!70!black] coordinates {(0,24.38)(1,31.13)(2,32.80)(3,32.90)(4,33.62)(5,34.21)(6,34.74)(7,35.85)(8,34.68)(9,36.68)(10,37.60)};
      \addlegendentry{ROUGE-L $\times$ 100}
    \end{axis}
  \end{tikzpicture}
  \caption{Few-shot accuracy of the frozen baseline against the number of demonstrations.}
  \label{fig:shots}
\end{figure}

Exact match rises from 15.0\% with no examples to 28.0\% with ten, and ROUGE-L from 0.244 to
0.376. Paired against the zero-shot answers on the same queries, ten demonstrations give
a gain of 13.0 points of EM (95\% bootstrap interval 8.0 to 18.5 points), with 29 queries
improved and 3 made worse. Most of the gain comes from the first two demonstrations:
they teach the output format, and the average answer length falls from 2.5 to about 1.5 words
(zero-shot answers are often full sentences). The cost grows linearly, since latency rises from 0.40\,s
to 2.23\,s and the prompt from 211 to 1,919 tokens. Table~\ref{tab:fsexamples} shows
that individual answers are unstable: for the hottest-temperature question the model
answered 400, 200, 300, 400, 100 and then 100 as $k$ increased, while the reference is
$>$500. Few-shot prompting therefore mostly fixes the answer format and does not
add reliable fire-specific perception.

\begin{table}[htbp]
  \centering
  \caption{Examples of frozen few-shot predictions for ten test queries (exact match marked $\checkmark$)}
  \label{tab:fsexamples}
  \small
  \begin{tabular}{P{2.0in}P{0.95in}P{1.0in}P{1.0in}}
    \toprule
    \textbf{Question (abridged)} & \textbf{Reference} & \textbf{$k{=}0$} & \textbf{$k{=}10$} \\
    \midrule
    Most intense hotspot location & Center & Upper right & Bottom-right \\
    Region of largest smoke plume & Center & Top-left & Left \\
    Fuel-bed continuity & Continuous & (full sentence) & Continuous $\checkmark$ \\
    Hottest temperature (\textdegree C) & $>$500 & 400 & 100 \\
    Safety risk of UAV position & Medium risk & High & High \\
    Scene share above 200\textdegree C & $<$5\% & 50\% & 0\% \\
    Region of largest smoke plume & Bottom-left & Top-left & Top-left \\
    Number of emergency vehicles & 0 & 0 $\checkmark$ & 0 $\checkmark$ \\
    Primary ground fuel type & Grass & Vegetation & Grass $\checkmark$ \\
    Buildings visible & Yes & No & No \\
    \bottomrule
  \end{tabular}
\end{table}

\subsection{QLoRA Fine-Tuning and Adapter Capacity}
Table~\ref{tab:ablation} lists every Qwen2-1.5B run that finished, including the extra ranks.
Fine-tuning lifts exact match from at most 28\% to between 84\% and 87\%.

\begin{table}[htbp]
  \centering
  \caption{All completed Qwen2-1.5B fine-tuning runs (2,482 test rows; best checkpoint by validation loss)}
  \label{tab:ablation}
  \footnotesize
  \setlength{\tabcolsep}{3pt}
  \resizebox{\textwidth}{!}{%
  \begin{tabular}{lrccccc}
    \toprule
    \textbf{Adapters, rank} & \textbf{Trainable params} & \textbf{EM (\%)} & \textbf{ROUGE-L} & \textbf{Best val.\ loss} & \textbf{Epochs (best)} & \textbf{Time (min)} \\
    \midrule
    $q,k,v,o$, $r{=}4$ & 5,816,320 & \textbf{86.62} & 0.8979 & 0.1236 & 14 (11) & 568 \\
    $q,k,v,o$, $r{=}8$ & 6,905,856 & 86.30 & 0.8928 & 0.1159 & 13 (10) & 512 \\
    $q,v$, $r{=}8$ & 5,816,320 & 86.14 & 0.8932 & 0.1209 & 12 (9) & 478 \\
    $q,v$, $r{=}16$ & 6,905,856 & 86.06 & 0.8937 & 0.1177 & 12 (9) & 456 \\
    all-linear, $r{=}8$ & 13,959,168 & 84.85 & 0.8858 & 0.1263 & 6 (3) & 283 \\
    $q,k,v,o$, $r{=}16$ & 9,084,928 & 84.13 & 0.8773 & 0.1270 & 10 (7) & 402 \\
    all-linear, $r{=}16$ & 23,191,552$^{*}$ & \todo{not run} & \todo{not run} & \todo{not run} & \todo{not run} & \todo{not run} \\
    \bottomrule
  \end{tabular}}\\[2pt]
  {\footnotesize Time is the total over all epochs on one T4 GPU. $^{*}$Computed, not executed. A rank-32 run ($q,k,v,o$) has no saved results.}
\end{table}

The four attention-only runs with ranks up to 16 on $q,v$ and ranks 4--8 on $q,k,v,o$ lie within 0.6
points of one another (86.06 to 86.62\%). With 2,482 test rows from only 73 images and one seed, differences
of this size are within sampling noise, so they should not be ranked. The two
configurations with the most adapter parameters, $q,k,v,o$ at $r{=}16$ and all-linear at $r{=}8$, are 1.3 to 2.5
points lower, and they also stopped earlier (best epoch 7 and 3). The smallest configuration
with near-best accuracy is therefore $q,v$ or $q,k,v,o$ at a low rank, which trains 5.8 million
parameters of which 4.7 million are the projectors. Adapter capacity beyond this did not help:
extra adapter parameters add training cost and, in these runs, lower accuracy, which is consistent with
over-fitting of the 592 training images. The result for all-linear at $r{=}16$ is not available.

\subsection{Llama-3.2-1B Backbone}
The same pipeline was run with \texttt{meta-llama/Llama-3.2-1B-Instruct} as the language model
(Table~\ref{tab:llama}). All of these runs score higher than the matching Qwen2 runs, from
0.2 points higher for $q,v$ at $r{=}8$ to 4.0 points higher for $q,k,v,o$ at $r{=}16$,
and the best result is 88.48\% EM for $q,k,v,o$ at $r{=}8$. The training settings differ between the Llama runs and from the Qwen2 runs
(e.g.\ learning rate $1.5\times10^{-4}$ versus $2\times10^{-4}$), so the backbone comparison is not
fully controlled. The ordering of the adapter variants is similar: the all-linear adapters did not beat
the attention-only ones, and rank 16 did not beat rank 8.

\begin{table}[htbp]
  \centering
  \caption{Completed Llama-3.2-1B runs (same data, split and metric)}
  \label{tab:llama}
  \footnotesize
  \setlength{\tabcolsep}{3pt}
  \resizebox{\textwidth}{!}{%
  \begin{tabular}{lrccccc}
    \toprule
    \textbf{Adapters, rank} & \textbf{Trainable params} & \textbf{EM (\%)} & \textbf{ROUGE-L} & \textbf{Best val.\ loss} & \textbf{Epochs} & \textbf{Time (min)} \\
    \midrule
    $q,k,v,o$, $r{=}8$ & 7,481,344 & \textbf{88.48} & 0.9121 & 0.1057 & 9 & 268 \\
    $q,k,v,o$, $r{=}16$ & 9,185,280 & 88.15 & 0.9073 & 0.1017 & 8 & 247 \\
    $q,v$, $r{=}16$ & 7,481,344 & 87.63 & 0.9043 & 0.1099 & 7 & 202 \\
    all-linear, $r{=}8$ & 11,413,504 & 87.63 & 0.9062 & 0.1096 & 9 & 310 \\
    $q,v$, $r{=}8$ & 6,629,376 & 86.38 & 0.8911 & 0.1233 & 10 & 302 \\
    \bottomrule
  \end{tabular}}
\end{table}

\begin{figure}[htbp]
  \centering
  \begin{tikzpicture}
    \begin{axis}[width=0.85\textwidth, height=6.4cm, xlabel={Trainable parameters (millions)},
      ylabel={Exact match (\%)}, xmin=5, xmax=15, ymin=83, ymax=89.5, grid=major,
      grid style={gray!25}, legend pos=north east, legend cell align=left]
      \addplot[only marks, mark=*, mark size=2.6pt, color=blue!70!black]
        coordinates {(5.82,86.14)(5.82,86.62)(6.91,86.06)(6.91,86.30)(9.08,84.13)(13.96,84.85)};
      \addlegendentry{Qwen2-1.5B}
      \addplot[only marks, mark=square*, mark size=2.6pt, color=red!70!black]
        coordinates {(6.63,86.38)(7.48,87.63)(7.48,88.48)(9.19,88.15)(11.41,87.63)};
      \addlegendentry{Llama-3.2-1B}
    \end{axis}
  \end{tikzpicture}
  \caption{Exact match against trainable parameters (adapters, projectors and modality embeddings) for all completed fine-tuning runs.}
  \label{fig:pareto}
\end{figure}

\subsection{Per-Category Results}
Table~\ref{tab:category} gives exact match by question category for the two runs whose notebooks
contain this breakdown. Presence and detection (mostly yes/no questions) is easiest and localization and direction
is hardest in both runs, which agrees with the larger number of answer choices
for the localization questions (about 6.5 on average, against about 2.4 for presence questions
\cite{wildfirevqa}).

\begin{table}[htbp]
  \centering
  \caption{Exact match (\%) by question category}
  \label{tab:category}
  \small
  \begin{tabular}{P{2.2in}cc}
    \toprule
    \textbf{Category} & \textbf{Qwen2, $q,k,v,o$ $r{=}16$} & \textbf{Llama, $q,v$ $r{=}8$} \\
    \midrule
    Presence and detection & 90.75 & 94.35 \\
    Distribution and segmentation & 85.96 & 90.41 \\
    Classification & 85.84 & 81.96 \\
    Flight planning & 85.27 & 85.27 \\
    Cross-modal reasoning & 76.37 & 81.85 \\
    Localization and direction & 71.23 & 74.66 \\
    \bottomrule
  \end{tabular}
\end{table}

\section{Analysis}
Table~\ref{tab:results} shows the effect of fine-tuning and of adapter capacity. Four points stand out.
First, task-specific training matters far more than prompting: the best frozen few-shot result is 28\%,
while every fine-tuned model scores above 84\%. This gap should be read with the caveats of Section~7.1
(different model size, test set and input handling), but it is large enough to be
meaningful. Second, a small adapter budget is enough. The best Qwen2 run uses $q,k,v,o$ at rank 4, and
the all-linear and rank-16 attention adapters, which add 3 to 8 million trainable parameters over the smallest
runs, do not improve accuracy. Third, the choice of language backbone matters at least as much as the
adapter configuration in these experiments: Llama-3.2-1B gave 0.2 to 4.0 points more than Qwen2-1.5B for matching
adapters. Fourth, the hard categories are those with many answer options, such as localization and
direction and cross-modal reasoning, where one error type is confusing neighbouring image regions or
temperature bins.

\section{Limitations}
(i) Each configuration was trained once with a single seed, and the test set contains only 73 image pairs,
so differences of about one point are not reliable. (ii) The data come from one burn site (Sycan Marsh); generalisation to other
sites was not tested. (iii) The few-shot baseline and the fine-tuned models use different language models, test
sets and image-path handling. (iv) Answers are scored by exact string match without the answer options
in the prompt, so a correct answer worded differently counts as wrong. (v) The all-linear $r{=}16$ and
$q,k,v,o$ $r{=}32$ Qwen2 notebooks have no saved results. (vi) The Qwen2 language component is paired with a CLIP
encoder and randomly initialised projectors rather than its own vision tower, so the model is not a native Qwen2-VL.

% ---------------- CHAPTER 8 ------------------------------------------
\clearpage
\chapter{Conclusion and Future Scope}

\section{Conclusion}
FireVLM shows that a compact vision-language model, built from a frozen CLIP encoder, trainable projectors and a
4-bit Qwen2-1.5B language model with QLoRA adapters, can answer wildfire questions about paired RGB and thermal
imagery on free-tier T4 GPUs. It reaches 84.1--86.6\% exact match on 2,482 held-out WildFireVQA questions
while training 5.8--14.0 million parameters, compared with 28\% for a frozen Qwen2.5-VL-3B model given ten
examples. Comparing adapter placement and rank shows that low-rank attention adapters give most of the gain:
the best Qwen2 configuration ($q,k,v,o$, rank 4) scores 86.62\%, and larger adapters did not improve it. With a
Llama-3.2-1B backbone the same pipeline reached 88.48\% ($q,k,v,o$, rank 8).

\section{Future Scope}
\begin{itemize}[leftmargin=0.5in, itemsep=2pt]
  \item Run the missing all-linear $r{=}16$ and $q,k,v,o$ $r{=}32$ configurations, repeat runs over several seeds, and evaluate the frozen baseline on the same 2,482-row test split.
  \item Evaluate on other fire sites (Willamette and Shoetank) and use the radiometric thermal TIFF data to test generalisation.
  \item Explore other adapter methods, larger CLIP-style or native Qwen2-VL vision towers, and fine-tuning the vision encoder.
  \item Score answers with the multiple-choice protocol of WildFireVQA so results can be compared with published baselines.
  \item Add temporal reasoning over video sequences of a fire.
  \item Deploy on an edge device for on-board drone inference.
\end{itemize}

% =====================================================================
%  REFERENCES
% =====================================================================
\clearpage\phantomsection
\addcontentsline{toc}{chapter}{REFERENCES}
\markboth{References}{}
\begin{singlespace}
\begin{thebibliography}{99}

\bibitem{antol2015}
S. Antol \textit{et al.}, ``VQA: Visual Question Answering'', Proceedings of the
IEEE International Conference on Computer Vision (ICCV), pp. 2425-2433, 2015.

\bibitem{bai2023}
J. Bai \textit{et al.}, ``Qwen-VL: A Versatile Vision-Language Model for
Understanding, Localization, Text Reading, and Beyond'', arXiv:2308.12966, 2023.

\bibitem{wang2024}
P. Wang \textit{et al.}, ``Qwen2-VL: Enhancing Vision-Language Model's Perception
of the World at Any Resolution'', arXiv:2409.12191, 2024.

\bibitem{bai2025}
S. Bai \textit{et al.}, ``Qwen2.5-VL Technical Report'', arXiv:2502.13923, 2025.

\bibitem{yang2024}
A. Yang \textit{et al.}, ``Qwen2 Technical Report'', arXiv:2407.10671, 2024.

\bibitem{mithan}
MIT HAN Lab, ``Qwen2-VL-1.5B-Instruct'', Hugging Face model repository, \url{https://huggingface.co/mit-han-lab/Qwen2-VL-1.5B-Instruct}.

\bibitem{llama32}
Meta AI, ``Llama-3.2-1B-Instruct'', Hugging Face model repository, \url{https://huggingface.co/meta-llama/Llama-3.2-1B-Instruct}.

\bibitem{radford2021}
A. Radford \textit{et al.}, ``Learning Transferable Visual Models From Natural
Language Supervision'', Proceedings of the 38th International Conference on Machine
Learning (ICML), PMLR 139, pp. 8748-8763, 2021.

\bibitem{brown2020}
T. Brown \textit{et al.}, ``Language Models are Few-Shot Learners'', Advances in
Neural Information Processing Systems (NeurIPS), Vol. 33, pp. 1877-1901, 2020.

\bibitem{hu2022}
E. J. Hu \textit{et al.}, ``LoRA: Low-Rank Adaptation of Large Language Models'',
International Conference on Learning Representations (ICLR), 2022.

\bibitem{dettmers2023}
T. Dettmers, A. Pagnoni, A. Holtzman, L. Zettlemoyer, ``QLoRA: Efficient
Finetuning of Quantized LLMs'', Advances in Neural Information Processing
Systems (NeurIPS), Vol. 36, 2023.

\bibitem{flame3}
B. Hopkins, L. ONeill, M. Marinaccio, E. Rowell, R. Parsons, S. Flanary, I. Nazim,
C. Seielstad, F. Afghah, ``FLAME 3 Dataset: Unleashing the Power of Radiometric Thermal
UAV Imagery for Wildfire Management'', arXiv:2412.02831, 2024.

\bibitem{wildfirevqa}
M. Habibpour, N. Alipour Talemi, J. Spodnik, C. J. Khoury, F. Afghah, ``WildFireVQA: A
Large-Scale Radiometric Thermal VQA Benchmark for Aerial Wildfire Monitoring'',
arXiv:2604.20190, 2026. Data and code: \url{https://github.com/mobiiin/WildFire_VQA}.

\end{thebibliography}
\end{singlespace}

\end{document}
